\chapter{The translator metric and graphical calibrations}\label{sec:translator-metric-calibration}
Ilmanen's metric identifies translators with minimal hypersurfaces for a different ambient metric. The source discussion is \cite{HIMW19}, pp.~9--10. We distinguish the two metrics and prove the graphical comparison by a closed weighted form. The form is defined throughout the graph cylinder, which specifies where competitors may lie.

\begin{proposition}[The conformal translator metric]\label{prop:translator-conformal-metric}
For fixed $v\in\mathbb R^{n+1}$ and $n\ge1$, put $h(x)=\langle x,v\rangle$ and $\bar g_v=e^{2h/n}\bar g_{\rm Eucl}$. On every immersed $n$-hypersurface, the $\bar g_v$ area measure is $e^h dA$. Consequently its translator equation is equivalent to $\bar g_v$ minimality. If $v\ne0$, this ambient metric is incomplete.
\end{proposition}
\begin{proof}
The induced metric is multiplied by $e^{2h/n}$ in each of $n$ tangent directions; its volume density is multiplied by $e^h$. Thus equality of the area functionals and Proposition \ref{prop:translator-weighted-area-example} give equivalence of stationarity. Along the ray $x(t)=-tv/|v|$, $t\ge0$, the metric length is
\[
 \int_0^\infty e^{-|v|t/n}\,dt=\frac n{|v|}<\infty.
\]
Points going to infinity along this ray form a metric Cauchy sequence with no limit in Euclidean space, whose usual local topology agrees with that of the smooth positive metric. Hence the metric is incomplete. The calculation does not assert incompleteness of every induced translator metric.
\end{proof}

\begin{proposition}[Normal speeds in the two metrics]\label{prop:translator-metric-normal-speed}
If $\nu$ is a Euclidean unit normal, then $\bar\nu=e^{-h/n}\nu$ is a $\bar g_v$ unit normal. A variation with Euclidean normal speed $f$ has $\bar g_v$ normal speed $\bar f=e^{h/n}f$. In particular the normal component of an ambient translation is $e^{h/n}\langle a,\nu\rangle$ in the conformal metric.
\end{proposition}
\begin{proof}
Conformal scaling preserves orthogonality and multiplies vector lengths by $e^{h/n}$. Hence $f\nu=e^{h/n}f\bar\nu$, proving the assertions. The Euclidean angle solves the weighted equation of Proposition \ref{prop:translator-angle-jacobi}; the conformal normal speed solves the corresponding minimal Jacobi equation. These are related formulations, not the same scalar operator applied to the same normal-speed function.
\end{proof}

\begin{proposition}[A closed form calibrates a translating graph]\label{prop:translator-graph-weighted-calibration}
Suppose $u\in C^\infty(\Omega)$, with $\Omega\subset\mathbb R^n$ open, solves the graph translator equation with velocity $ce_{n+1}$. On $U=\Omega\times\mathbb R$, set
\[
 N(x,z)=\frac{(-Du(x),1)}{\sqrt{1+|Du(x)|^2}},\qquad
 \rho(x,z)=e^{cz},\qquad \omega=\iota_{\rho N}(dx_1\wedge\cdots\wedge dx_n\wedge dz).
\]
Then $d\omega=0$, the Euclidean comass of $\omega$ is $\rho$, and its restriction to the graph with the orientation induced by $N$ is its weighted area form. Equivalently $\omega$ has comass one for the conformal metric of Proposition \ref{prop:translator-conformal-metric}.
\end{proposition}
\begin{proof}
Put $W=\sqrt{1+|Du|^2}$. Because $N$ is independent of $z$,
\[
 \operatorname{div}_{\mathbb R^{n+1}}(\rho N)
 =\rho\left(-\operatorname{div}_{\mathbb R^n}\frac{Du}{W}+\frac cW\right)=0
\]
by Proposition \ref{prop:translator-graph-example}. The exterior derivative of the contraction with ambient volume is this divergence times ambient volume, so the form is closed. Its Euclidean comass is the length of $\rho N$, namely $\rho$, with equality on the oriented plane normal to $N$. For the conformal metric an orthonormal $n$-frame is the Euclidean one multiplied by $e^{-cz/n}$, canceling that factor. This proves the claimed calibration without requiring a global minimal-foliation theorem.
\end{proof}

\begin{theorem}[Weighted minimization within the graph cylinder]\label{thm:translator-graph-current-minimization}
Let $T$ be the current of integration over a compact smooth piece of the oriented translating graph in Proposition \ref{prop:translator-graph-weighted-calibration}. If $S$ is a compactly supported integer rectifiable $n$-current of finite mass in $U$, and $S-T=\partial R$ for a compactly supported $(n+1)$-current $R$ in $U$, then
\[
 \int\rho\,d\|S\|\ge\int\rho\,d\|T\|.
\]
Equality holds exactly when the oriented approximate tangent of $S$ is calibrated by $\omega$ almost everywhere. This gives a current formulation of the graphical minimizing assertion in \cite{HIMW19}, pp.~9--10, with its ambient-domain condition explicit.
\end{theorem}
\begin{proof}
Choose a cutoff equal to one near all compact supports. Since $d\omega=0$ and its derivative vanishes near the support of $R$ after cutoff, Stokes' identity for currents gives $S(\omega)=T(\omega)$. The latter equals the weighted mass of $T$. The former is bounded above by the weighted mass of $S$, because its oriented rectifiable representation pairs a unit simple tangent vector with a form of comass $\rho$. The difference is the integral of a nonnegative calibration defect; it vanishes exactly when that defect is zero almost everywhere. The argument is the weighted version of Theorem \ref{thm:calibrated-current-minimizes}. It compares homology classes inside $U$, not arbitrary competitors that leave the graph cylinder.
\end{proof}

\begin{corollary}[Strict minimizing comparison between graphs with the same boundary]\label{cor:translator-fixed-boundary-graph-minimum}
Let $D\Subset\Omega$ be a connected smooth bounded domain and let $u$ be as in Proposition \ref{prop:translator-graph-weighted-calibration}. For every smooth $w$ on $\overline D$ with $w=u$ on $\partial D$,
\[
 \int_D e^{cw}\sqrt{1+|Dw|^2}\,dx
 \ge\int_D e^{cu}\sqrt{1+|Du|^2}\,dx.
\]
Equality holds if and only if $w=u$. In particular there is at most one smooth solution of this translator Dirichlet problem. Existence is a separate problem.
\end{corollary}
\begin{proof}
The vertical homotopy $(x,t)\mapsto(x,(1-t)u(x)+tw(x))$ supplies a compact filling of the difference of the two graph currents in $\Omega\times\mathbb R$. Its side term is zero because the functions agree on $\partial D$. The preceding theorem gives the inequality. Equality makes the upward unit normals agree everywhere, first almost everywhere by the defect formula and then everywhere by continuity. Their last components are positive, so equality of normals gives $Dw=Du$. Connectedness makes $w-u$ constant, and the boundary condition makes the constant zero.

If two solutions are smooth on $\overline D$ and have the same boundary data, apply this calibration calculation in both directions. No extension as a solution beyond $\partial D$ is needed: pull the closed form back to the vertical homotopy on $\overline D\times[0,1]$ and use Stokes there. It is closed on the interior, smooth up to the boundary, and the side term vanishes because the two heights agree. Thus each graph has weighted area at most the other's. Their areas are equal, and the equality statement gives uniqueness. No joint convexity of the integrand in height and gradient is used.
\end{proof}

\begin{proposition}[A local area bound for a translating graph]\label{prop:translator-graph-local-area-bound}
Let $M$ be the graph in Proposition \ref{prop:translator-graph-weighted-calibration}. If $K\Subset\Omega\times\mathbb R$ is a smooth bounded domain whose boundary meets $M$ transversely, then
\[
 2\int_{M\cap K}e^{cz}\,dA\le\int_{\partial K}e^{cz}\,dA.
\]
In particular, for a transverse Euclidean ball $B_r(p)\Subset\Omega\times\mathbb R$,
\[
 \mathcal H^n(M\cap B_r(p))
 \le\tfrac12 e^{2|c|r}\mathcal H^n(\partial B_r(p)).
\]
The estimate alone is a mass bound, not a theorem of smooth convergence for arbitrary dimensions or boundary configurations.
\end{proposition}
\begin{proof}
The graph separates the cylinder into $U^+=\{z>u(x)\}$ and $U^-=\{z<u(x)\}$. Apply the divergence theorem to the divergence-free field $Y=e^{cz}N$ on $K\cap U^+$ and $K\cap U^-$, whose boundaries are piecewise smooth by transversality. The outward normal on the graph is respectively $-N$ and $N$. Therefore
\[
 \int_{M\cap K}\rho=\int_{\partial K\cap U^+}\langle Y,\eta\rangle,
 \qquad
 \int_{M\cap K}\rho=-\int_{\partial K\cap U^-}\langle Y,\eta\rangle.
\]
Since $|Y|=\rho$, summing proves the weighted bound. On $B_r(p)$ the maximum of $e^{cz}$ divided by its minimum is at most $e^{2|c|r}$. Use the minimum on the interior integral and the maximum on the boundary integral to obtain the unweighted bound. These computations give the area comparison described in \cite{HIMW19}, p.~10, with the domain and transversality hypotheses specified.
\end{proof}
